% Theorem-facing replacement for Section 3.

\section{Continuous profile roots}
\label{continuous-profile-roots}

For a positive class size $u$, put
\[
  d_u:=2^{\binom u2}u!.
\]
Write $u=\alpha-i$, so $i$ is the deficit from the phase center $\alpha$.
We use
\begin{equation*}
  S_+:=\{-1,0,1,2,\ldots\},
  \qquad
  S_4:=\{2,3,4,5\}.
  \tag{3.1}
\end{equation*}
At finite $n$, the unrestricted support is
$S_+^{(n)}:=\{-1,0,\ldots,\alpha-1\}$, because all class sizes are
positive. Whenever a finite-$n$ sum is indexed by $S_+$, this cutoff is
understood.

For a real class count $k>0$, define
\begin{equation*}
  L_{\mathbf k}
  :=n\log n-n+k-\sum_i k_i\log(k_i d_{\alpha-i}),
  \tag{3.2}
\end{equation*}
with the convention $0\log0=0$, and let $L_S(n,k)$ be its maximum over
real $k_i\ge0$ satisfying
\begin{equation*}
  \sum_{i\in S}k_i=k,
  \qquad
  \sum_{i\in S}(\alpha-i)k_i=n.
  \tag{3.3}
\end{equation*}
For a feasible $k$, set
\begin{equation*}
  s:=\frac nk,
  \qquad
  T:=\alpha-s=\alpha-\frac nk.
  \tag{3.4}
\end{equation*}
Thus $T$ is the mean deficit of the normalized profile.

\begin{lemma}[Uniform root, slope, and finite-dual package]
Let
\begin{equation*}
  s_0:=\alpha_0-1-\frac2{\log 2},
  \qquad
  T_0:=\alpha-s_0
      =1+\frac2{\log 2}-\delta.
  \tag{3.5}
\end{equation*}
There is a fixed $A_*>0$ such that, for
$S\in\{S_+,S_4\}$ and $0\le c\le\log 2$, the equation
\[
  L_S(n,k)+ck=0
\]
has, for all sufficiently large $n$, a unique zero $r_{S,c}$ in the
corridor
\[
  \left|\frac nk-s_0\right|
  \le A_*\frac{\log\log n}{\log n}.
\]
Uniformly in $S$, $c$, and the phase,
\begin{equation*}
  \frac n{r_{S,c}}
  =s_0+O\!\left(\frac{\log\log n}{\log n}\right),
  \qquad
  \alpha-\frac n{r_{S,c}}
  =T_0+O\!\left(\frac{\log\log n}{\log n}\right).
  \tag{3.6}
\end{equation*}

For every fixed $A>0$, uniformly when
$|n/k-s_0|\le A\log\log n/\log n$,
\begin{equation*}
  \frac{\partial}{\partial k}\{L_S(n,k)+ck\}
  =\frac2{\log 2}(\log n)^2
   +O_A((\log n)(\log\log n)).
  \tag{3.7}
\end{equation*}
Equivalently, the deterministic normalized slope error
\[
  \varepsilon_{n,A}^{\mathrm{slope}}
  :=
  \max_{\substack{S\in\{S_+,S_4\}\\0\le c\le\log 2}}
  \sup_{|n/k-s_0|\le A\log\log n/\log n}
  \left|
    \frac{1}{(\log n)^2}
    \frac{\partial}{\partial k}\{L_S(n,k)+ck\}
    -\frac2{\log 2}
  \right|
\]
satisfies $\varepsilon_{n,A}^{\mathrm{slope}}\to0$.

For two supports $R,S\in\{S_+,S_4\}$ at the same feasible $k$,
\begin{equation*}
  \frac{L_R(n,k)-L_S(n,k)}{k}
  =\mathcal F_{n,R}(T)-\mathcal F_{n,S}(T),
  \tag{3.8}
\end{equation*}
where
\begin{equation*}
  \mathcal F_{n,S}(T)
  :=
  \max_{\substack{\sum_i p_i=1\\\sum_i i p_i=T}}
  \left[
    -\sum_i p_i\log p_i+\sum_i p_i h_n(i)
  \right].
  \tag{3.8a}
\end{equation*}
Fix the compact target interval
\begin{equation*}
  K_*:=
  \left[
    \frac2{\log 2}-\frac1{10},
    1+\frac2{\log 2}+\frac1{10}
  \right]
  \subset(2,5).
  \tag{3.9a}
\end{equation*}
Then, uniformly for $T\in K_*$,
\begin{equation*}
  \mathcal F_{n,S}(T)
  \longrightarrow
  \mathcal F_S(T)
  :=
  \max_{\substack{\sum_i p_i=1\\\sum_i i p_i=T}}
  \left[
    -\sum_i p_i\log p_i
    -\frac{\log 2}{2}\sum_i i^2p_i
  \right].
  \tag{3.9}
\end{equation*}
If $\lambda_{n,S}(T)$ and $\lambda_S(T)$ are the finite and limiting
effective tilts, then there is a fixed $\Lambda>0$ such that
\[
  |\lambda_{n,S}(T)|+|\lambda_S(T)|\le2\Lambda
\]
for all sufficiently large $n$, and the deterministic dual error
\[
\begin{split}
  \varepsilon_n^{\mathrm{dual}}
  :=\max_{S\in\{S_+,S_4\}}
  \sup_{T\in K_*}
  \bigl(&|\mathcal F_{n,S}(T)-\mathcal F_S(T)|\\
        &+|\lambda_{n,S}(T)-\lambda_S(T)|\bigr)
\end{split}
\]
satisfies $\varepsilon_n^{\mathrm{dual}}\to0$.

For $S_4$, if $p_{n,i}(T)$ and $p_i(T)$ denote the finite and limiting
optimizers, then
\begin{equation*}
  \sup_{T\in K_*}\max_{2\le i\le5}
  |p_{n,i}(T)-p_i(T)|\longrightarrow0,
  \qquad
  \inf_{\substack{n\ \mathrm{large}\\T\in K_*}}
  \min_{2\le i\le5}p_{n,i}(T)>0.
  \tag{3.9b}
\end{equation*}
All error sequences and eventuality thresholds above are independent of
$\delta$.
\end{lemma}

\begin{proof}
Put $q:=\log 2$, $L:=\log n$, and $\ell:=\log\log n$.
Dividing (3.2) by $k=n/s$ and writing $p_i=k_i/k$ gives
\[
  \frac{L_S(n,n/s)}{n/s}
  =(s-1)L-s+1+\log s
   +\max_p\left[-\sum_i p_i\log p_i
                 -\sum_i p_i\log d_{\alpha-i}\right].
\]
There is an exact affine-plus-curved decomposition
\begin{equation*}
  -\log d_{\alpha-i}=A_n+B_ni+h_n(i),
  \tag{3.11}
\end{equation*}
where
\begin{equation*}
  A_n:=-\log d_\alpha,
  \qquad
  B_n:=q\alpha-\frac q2+\log\alpha,
  \tag{3.12}
\end{equation*}
and, for $i\ge0$,
\begin{equation*}
  h_n(i)
  :=-\frac q2i^2
    +\sum_{r=0}^{i-1}\log\!\left(1-\frac r\alpha\right).
  \tag{3.13}
\end{equation*}
For $i=-1$, direct subtraction gives
\[
  h_n(-1)=-\frac q2+\log\!\left(\frac{\alpha}{\alpha+1}\right).
\]
Thus
\begin{equation*}
  h_n(i)=-\frac q2i^2+O(i^2/\alpha)
  \quad\text{for every fixed }i,
  \qquad
  h_n(i)\le-\frac q2i^2
  \quad(i\ge-1).
  \tag{3.14}
\end{equation*}
The first estimate is uniform on every fixed finite set of deficits; the
second is global on the finite support.

\displayheading{Finite duals and optimizing tilts}
For $S\in\{S_+,S_4\}$, define
\[
  Z_{n,S}(\lambda)
  :=\sum_{i\in S}
    \exp\{\lambda i+h_n(i)\},
  \qquad
  M_{n,S}(\lambda)
  :=\frac{d}{d\lambda}\log Z_{n,S}(\lambda),
\]
with the cutoff $S_+^{(n)}$ in the finite unrestricted sum. Define the
limiting functions by
\[
  Z_S(\lambda)
  :=\sum_{i\in S}
    \exp\!\left(\lambda i-\frac q2i^2\right),
  \qquad
  M_S(\lambda):=\frac{d}{d\lambda}\log Z_S(\lambda).
\]
The target center satisfies
\begin{equation*}
  \frac2q\le T_0\le1+\frac2q,
  \tag{3.15}
\end{equation*}
so it lies in the interior of $K_*$. The same is true of every target in
$K_*$ relative to both supports.

Choose $\Lambda>0$ so that
\[
  M_S(-\Lambda)<\min K_*<\max K_*<M_S(\Lambda)
  \qquad(S=S_+,S_4).
\]
Such a choice exists because the limiting mean maps are strictly increasing
and their endpoint limits bracket $K_*$. For $m=0,1,2$ and
$|\lambda|\le\Lambda$, (3.14) gives the summable majorant
\[
  |i|^m\exp\!\left(\Lambda|i|-\frac q2i^2\right).
\]
On every fixed finite set of indices, the weights converge by (3.14). The
majorant makes the remaining Gaussian tail uniformly small, and the omitted
tail beyond the finite cutoff $\alpha-1$ is bounded by the same series.
Consequently the zeroth, first, and second tilted moments converge uniformly
on $[-\Lambda,\Lambda]$. In particular,
\[
  Z_{n,S}\to Z_S,
  \qquad
  M_{n,S}\to M_S,
  \qquad
  M'_{n,S}\to M'_S
\]
uniformly there.

Now
$M'_S(\lambda)=\operatorname{Var}_{S,\lambda}(i)>0$. By continuity and
compactness,
\[
  v_*:=\frac12
  \min_{\substack{S\in\{S_+,S_4\}\\|\lambda|\le\Lambda}}
  M'_S(\lambda)>0.
\]
For all sufficiently large $n$, one has
$M'_{n,S}(\lambda)\ge v_*$ on the same interval, and the finite mean maps
also bracket $K_*$. Hence, for every $T\in K_*$, the equations
\[
  M_{n,S}(\lambda_{n,S}(T))=T,
  \qquad
  M_S(\lambda_S(T))=T
\]
have unique solutions in $[-\Lambda,\Lambda]$. The inverse mean maps are
uniformly $v_*^{-1}$-Lipschitz, and therefore
\[
  \sup_{T\in K_*}
  |\lambda_{n,S}(T)-\lambda_S(T)|
  \le
  v_*^{-1}\sup_{|\lambda|\le\Lambda}
  |M_{n,S}(\lambda)-M_S(\lambda)|
  \longrightarrow0.
\]

The entropy functional in (3.8a) is strictly concave on its feasible
simplex. Its unique optimizer is
\[
  p_{n,S,T}(i)
  =\frac{
    \exp\{\lambda_{n,S}(T)i+h_n(i)\}
  }{Z_{n,S}(\lambda_{n,S}(T))},
\]
and the dual value is exactly
\[
  \mathcal F_{n,S}(T)
  =\log Z_{n,S}(\lambda_{n,S}(T))
   -\lambda_{n,S}(T)T.
\]
The corresponding formulas hold in the limit. Uniform convergence of the
partition functions and tilts proves (3.9) and
$\varepsilon_n^{\mathrm{dual}}\to0$. On the finite support $S_4$, the optimizer
formula also gives uniform coordinatewise convergence. The limiting
coordinates are positive and continuous in $T$ on the compact set $K_*$;
this proves (3.9b).

\displayheading{Exact cancellation between supports}
Substituting (3.11) into the normalized exponent gives the exact identity
\begin{equation*}
\begin{split}
  \Psi_{n,S}(s)
  &:=\frac{L_S(n,n/s)}{n/s}\\
  &=(s-1)L-s+1+\log s
    +A_n+B_n(\alpha-s)
    +\mathcal F_{n,S}(\alpha-s).
\end{split}
  \tag{3.10}
\end{equation*}
At a fixed feasible $k$, the values of $s$ and $T=\alpha-s$ are the same for
both supports. Every term in (3.10) except the final dual value is therefore
common, which proves the exact difference identity (3.8). No limiting
replacement is used in this cancellation.

\displayheading{The value and derivative at the phase center}
At $s=s_0$, let $T_0=\alpha-s_0$. Applying Stirling once to $A_n$ in
(3.10), and then simplifying the affine terms exactly, gives the scalar
expression
\begin{equation*}
\begin{split}
  &s_0\left(
    L-\log\alpha-\frac q2s_0+\frac q2
  \right)-L+T_0+1+\log s_0+\frac q2T_0^2\\
  &\hspace{42mm}
  -\frac12\log(2\pi\alpha)+O(1/L).
\end{split}
  \tag{3.15a}
\end{equation*}
Equation (2.1) gives the exact identity
\[
  \frac{qs_0}{2}=L-\ell+\log q-q.
\]
Since $\alpha=s_0+T_0$ and $T_0$ is uniformly bounded,
\[
  \log\alpha
  =\log s_0+\frac{T_0}{s_0}+O(L^{-2}),
  \qquad
  \log s_0=\ell+q-\log q+O(\ell/L).
\]
Consequently
\[
\begin{split}
  s_0\left(
    L-\log\alpha-\frac q2s_0+\frac q2
  \right)
  &=-T_0+\frac q2s_0+O(\ell)\\
  &=-T_0+L-\ell+\log q-q+O(\ell).
\end{split}
\]
Substitution in (3.15a) cancels the terms $L$ and $T_0$ and leaves
$O(\ell)$ uniformly in the phase. The dual term in (3.10) is $O(1)$ uniformly
on $K_*$, by the Gaussian majorant above. Therefore
\begin{equation*}
  \Psi_{n,S}(s_0)=O(\ell)
  \tag{3.16}
\end{equation*}
uniformly for both supports.

The phase formula for $\alpha$ and the expansion of $\log\alpha$ also give
\begin{equation*}
  B_n=2L-\ell+O(1)
  \tag{3.17}
\end{equation*}
uniformly in $\delta$. Strict concavity and the optimizer formula imply
\[
  \frac{d}{dT}\mathcal F_{n,S}(T)=-\lambda_{n,S}(T).
\]
Differentiating (3.10), with $T=\alpha-s$, yields
\begin{equation*}
  \Psi'_{n,S}(s)
  =L-1+s^{-1}-B_n+\lambda_{n,S}(T)
  =-L+O_A(\ell)
  \tag{3.18}
\end{equation*}
uniformly whenever $|s-s_0|\le A\ell/L$. Indeed, this corridor lies inside
$K_*$ for all sufficiently large $n$, because $K_*$ has a fixed margin around
the complete range of $T_0$.

\displayheading{Existence, uniqueness, and slope of the roots}
Set
\[
  \Psi_{n,S,c}(s):=\Psi_{n,S}(s)+c.
\]
By (3.16), there is an absolute $C_0$ such that
$|\Psi_{n,S,c}(s_0)|\le C_0\ell$ for both supports and every
$0\le c\le q$. By (3.18), after one phase-independent eventuality threshold,
\[
  \Psi'_{n,S,c}(s)\le-\frac12L
\]
throughout each fixed corridor. Choose $A_*>2C_0$. Integration gives
\[
\begin{aligned}
  \Psi_{n,S,c}\!\left(s_0-A_*\frac\ell L\right)
  &\ge\left(\frac{A_*}{2}-C_0\right)\ell>0,\\
  \Psi_{n,S,c}\!\left(s_0+A_*\frac\ell L\right)
  &\le\left(C_0-\frac{A_*}{2}\right)\ell<0.
\end{aligned}
\]
The function is strictly decreasing in $s$ on this interval, so it has one and
only one zero there. Since $k=n/s>0$, this is equivalent to the unique root
$r_{S,c}$ of $L_S(n,k)+ck=0$, and proves (3.6).

Finally,
\[
  L_S(n,k)+ck=k\Psi_{n,S,c}(s),
  \qquad
  \frac{ds}{dk}=-\frac sk.
\]
Hence
\begin{equation*}
\begin{split}
  \frac{\partial}{\partial k}\{L_S(n,k)+ck\}
  &=\Psi_{n,S,c}(s)-s\Psi'_{n,S}(s)\\
  &=\frac2qL^2+O_A(L\ell),
\end{split}
  \tag{3.19}
\end{equation*}
because $s=(2/q)L+O_A(\ell)$, (3.18) holds, and
$\Psi_{n,S,c}(s)=O_A(\ell)$ throughout the corridor. This proves (3.7) and
$\varepsilon_{n,A}^{\mathrm{slope}}\to0$ with one eventuality threshold for
the complete phase.
\end{proof}